% Draft supplement for the next thesis build; not yet included in the released v2 PDF.
\subsection{Scope of stereo and accompaniment evidence}
Liu et al.'s M2S-ADD converts mono speech into synthesized stereo before
classification, with detector experiments on ASVspoof2019 LA
(\href{https://www.isca-archive.org/interspeech_2023/liu23v_interspeech.pdf}{Interspeech2023 primary paper},
doi:10.21437/Interspeech.2023-2335).
This result does not establish a native-stereo anomaly in generated music.
Inter-channel coherence, mid/side energy and spatial-width dynamics therefore
remain candidate production descriptors, not admitted authenticity cues.
Any later investigation must control native channel count, duplicated mono,
deliberate panning/widening, common gain, channel exchange and lossy encoding.

The \href{https://arxiv.org/abs/2604.16254v2}{ArtifactNet preprint, version2}
describes learned codec-residual extraction followed by harmonic/percussive
decomposition and classification. Its author-reported benchmark results are
not reproduced in this thesis. Residual extraction is treated as an alternative
pipeline, not automatically as an independent perceptual phenomenon.

Chen et al.'s \href{https://research.nvidia.com/labs/twn/publication/icassp_2026_singfakeinstrumental/}{ICASSP2026 author abstract}
reports primarily an augmentation role for accompaniment in singing-voice
deepfake detection rather than intrinsic rhythm or harmony evidence. This
motivates explicit feature ablations, but cannot by itself establish the same
conclusion for whole-song generation detection. The task definitions differ.
Only the abstract and version record of the latter two works were inspected
in this supplementary checkpoint; these notes are not complete methodological
replications. Further source and admission notes are retained in
\texttt{LITERATURE\_SCOPE\_CHECK\_20260907\_EN.md}.
